\documentclass{article}
\usepackage[margin=2.54cm]{geometry}
\usepackage{float}
\usepackage{amssymb}
\usepackage[notes,backend=biber]{biblatex-chicago}

\author{Anderson Lai}
\title{On the Optimal ``Minimum Point Value''}
\date{\today}

\addbibresource{references.bib}

\begin{document}
\maketitle

\begin{abstract}
    We aim to find the optimal ``minimum point value'' used in \verb|texnique.js| to maximize the amount of points earned in light of constraining factors.
\end{abstract}

\section{Introduction}
Initially, this problem may be trivial. Simply, set the ``minimum point value'' (MPV) to the highest available point so that only the highest point problem is solved. Whilst this appears to be a valid solution, it encounters some problems:
\begin{enumerate}
    \item Scores exceeding $2000$ points are automatically rejected by the leaderboard; they can not be inputted into the leaderboard.
    \item Pruning for only these high value problems results in having to skip many problems, which incurs a time penalty. Hence, fewer problems can be solved.
    \item The end screen renders all the skipped problems. Since so many problems are skipped, the end screen must render an absurd amount of problems, leading to the browser tab stalling and crashing.
\end{enumerate}
The main issue addressed here is the first issue. Whilst the second issue is somewhat relevant, the time penalty is minimal, and so it may be ignored. Moreover, as we will see later (or is immediately obvious), the third problem is naturally mitigated as we lower the MPV.

\section{Calculations}
It may be useful to have an estimate of the number of problems that the boat can solve in the timed mode (3 minutes). When setting the MPV to $13$ points, an average of $96.33$ problems can be solved. With the MPV at $17$ points, around $93$ problems can be solved (this reinforces that issue $2$ can be ignored). It is also useful to have the number of questions for each point value available\autocite{texnique-repo}.
\begin{table}[H]
    \begin{tabular}{|c|c|}
        \hline
        \textbf{Point Value} & \textbf{Question Count} \\
        \hline
        \hline
    
        $1$ & $2$ \\
        $2$ & $6$ \\
        $3$ & $9$ \\
        $4$ & $15$ \\
        $5$ & $22$ \\
        $6$ & $19$ \\
        $7$ & $16$ \\
        $8$ & $24$ \\
        $9$ & $14$ \\
        $10$ & $8$ \\
        $11$ & $8$ \\
        $12$ & $9$ \\
        $13$ & $11$ \\
        $14$ & $8$ \\
        $15$ & $5$ \\
        $16$ & $5$ \\
        $17$ & $1$ \\
        $21$ & $1$ \\
        $22$ & $1$ \\
        $33$ & $1$ \\

        \hline
    \end{tabular}
    \centering
    \caption{Question Count by Point Value}\label{table:question_count}
\end{table}

We can now calculate the \textit{expected value} of a problem's point value for a problem that has at least a minimum number of points. This is done by calculating the total number of questions available that meet this point threshold. Next, the relative abundance of each point value is calculated. Then, the standard expected value formula ($\mathbb{E}[X] = \sum_{i=1}^{N} p_i x_i$) is applied.
\begin{table}[H]
    \begin{tabular}{|c|c|}
        \hline
        \textbf{Minimum Point Value} & \textbf{Expected Value} \\
        \hline
        \hline
    
        $1$ & $8.29$ \\
        $2$ & $8.37$ \\
        $3$ & $8.59$ \\
        $4$ & $8.89$ \\
        $5$ & $9.37$ \\
        $6$ & $10.10$ \\
        $7$ & $10.79$ \\
        $8$ & $11.43$ \\
        $9$ & $12.57$ \\
        $10$ & $13.43$ \\
        $11$ & $13.98$ \\
        $12$ & $14.55$ \\
        $13$ & $15.24$ \\
        $14$ & $16.36$ \\
        $15$ & $17.71$ \\
        $16$ & $19.22$ \\
        $17$ & $23.25$ \\
        $21$ & $25.33$ \\
        $22$ & $27.50$ \\
        $33$ & $33.00$ \\

        \hline
    \end{tabular}
    \centering
    \caption{Expected Value for at least a Minimum Point Value}\label{table:expected_value}
\end{table}

As expected, the expected value increases as the MPV increases. However, this is the expected value \textit{per point}. We are instead interested in the expected value at the end of the 3 minutes for the timed mode. Hence, we must multiply each value by $95$ (an average number of points that can be solved in 3 minutes) to obtain our final expected value.

\begin{table}[H]
    \begin{tabular}{|c|c|}
        \hline
        \textbf{Minimum Point Value} & \textbf{Expected Value} \\
        \hline
        \hline

        $1$ & $787.7$ \\
        $2$ & $795.3$ \\
        $3$ & $815.8$ \\
        $4$ & $844.3$ \\
        $5$ & $889.8$ \\
        $6$ & $959.4$ \\
        $7$ & $1025.5$ \\
        $8$ & $1085.6$ \\
        $9$ & $1194.1$ \\
        $10$ & $1275.9$ \\
        $11$ & $1328.1$ \\
        $12$ & $1382.0$ \\
        $13$ & $1448.0$ \\
        $14$ & $1554.5$ \\
        $15$ & $1682.9$ \\
        $16$ & $1826.1$ \\
        $17$ & $2208.8$ \\
        $21$ & $2406.7$ \\
        $22$ & $2612.5$ \\
        $33$ & $3135.0$ \\

        \hline
    \end{tabular}
    \centering
    \caption{Expected Value After 3 Minutes}\label{table:final_expected_value}
\end{table}

\section{Conclusion}
Our desired MPV immediately becomes obvious: \textbf{17}. It meets the leaderboard cap of $2000$ points by a margin of $200$ points ($10\%$), whilst not exceeding the cap by too much, avoiding having to skip too many problems (addressing issues $2$ and $3$). This margin also provides a buffer to account for the ignoring of the time-penalty issue (issue $2$).

However, there is some criticism of this MPV. With the current implementation of the bot, pruning only for questions with points $\geq 17$ does have a small, but noticeable, time penalty. Mathematically, the number of available problems decreases from $9$ problems ($\approx 4.86\%$ of problems) to only $4$ problems ($\approx 2.16\%$ of problems)! 

Moreover, when pruning for the last few problems to reach exactly $2000$ points, the bot prunes for the highest point problem that does not exceed $2000$ points. This leads to an increase in the number of skipped problems, as the majority of the time, the $33$ point problem is being pruned for, and there is only $1$ of this type of problem. 

Nevertheless, whilst the end screen does experience meaningful lag with the MPV set to $17$, it is not enough to fully crash the browser. Hence, setting the MPV to $17$ is a valid and acceptable solution.

\newpage
\printbibliography

\end{document}